\documentclass[11pt]{article}
\usepackage[margin=1in]{geometry}
\usepackage[T1]{fontenc}
\usepackage[utf8]{inputenc}
\usepackage{lmodern}
\usepackage{hyperref}
\usepackage{xcolor}
\usepackage{fancyvrb}
\usepackage{listings}
\usepackage{CJKutf8}
\lstset{basicstyle=\ttfamily\small,columns=fullflexible,breaklines=true,upquote=true,frame=single,framerule=0.2pt,tabsize=2}

\title{\texttt{\detokenize{dataloader/kuairec.py}} (Q\&A Documentation)}
\author{}
\date{}
\begin{document}
\begin{CJK*}{UTF8}{gbsn}
\maketitle

\paragraph{Q: What does this file implement at a high level?}
\textbf{A:} This file defines the PyTorch Dataset/DataLoader wrapper for the preprocessed KuaiRec pickle. It converts pandas columns into tensors, casts dtypes based on the schema, and yields (features\textbackslash\{\}\_dict, label) batches.

\paragraph{Q: Why is dtype handling critical here?}
\textbf{A:} Embedding layers require integer indices, while continuous features and masks require float tensors. Wrong dtypes can crash forward passes or silently break training.

\paragraph{Q: What does this file export (public API)?}
\textbf{A:} The public API is the set of classes/functions other modules import (sometimes re-exported by package initializers). The key top-level definitions detected here are:

\begin{Verbatim}[fontsize=\small]
class KUAIRECDataset  # methods: __init__, format, __getitem__, __len__
class KUAIRECDataLoader  # methods: __init__, __getitem__
\end{Verbatim}

\paragraph{Q: What are the main dependencies (imports) and why do they matter?}
\textbf{A:} Imports show whether this file is model code (torch), data code (pandas), or evaluation/analysis (sklearn). They also determine runtime requirements.

\vspace{1mm}
\VerbatimInput[firstline=1,lastline=50,fontsize=\small]{/workspace/dataloader/kuairec.py}


\paragraph{Q: What is \texttt{\detokenize{KUAIRECDataset}} and what responsibility does it have?}
\textbf{A:} Load a KuaiRec Dataset

\paragraph{Q: What are the important methods on \texttt{\detokenize{KUAIRECDataset}}?}
\textbf{A:} In this repo, methods like forward/loss/predict define computation and training targets. The methods found are:

\begin{Verbatim}[fontsize=\small]
__init__
format
__getitem__
__len__
\end{Verbatim}

\paragraph{Q: What does the implementation look like for \texttt{\detokenize{KUAIRECDataset}}?}
\textbf{A:} Excerpt from the file:

\vspace{1mm}
\VerbatimInput[firstline=9,lastline=40,fontsize=\scriptsize]{/workspace/dataloader/kuairec.py}


\paragraph{Q: What is \texttt{\detokenize{KUAIRECDataLoader}} and what responsibility does it have?}
\textbf{A:} Load KUAIRANK16DataLoader for torch train/eval

:param dataset\_path: dataset path

\paragraph{Q: What are the important methods on \texttt{\detokenize{KUAIRECDataLoader}}?}
\textbf{A:} In this repo, methods like forward/loss/predict define computation and training targets. The methods found are:

\begin{Verbatim}[fontsize=\small]
__init__
__getitem__
\end{Verbatim}

\paragraph{Q: What does the implementation look like for \texttt{\detokenize{KUAIRECDataLoader}}?}
\textbf{A:} Excerpt from the file:

\vspace{1mm}
\VerbatimInput[firstline=43,lastline=66,fontsize=\scriptsize]{/workspace/dataloader/kuairec.py}


\paragraph{Q: Example: end-to-end data flow involving this file}
\textbf{A:} Core pipeline shape (schematic):

\begin{Verbatim}[fontsize=\small]
dataloaders = KUAIRECDataLoader("kuairec", "./dataset/kuairec/kuairec_data.pkl", device, bsz=2048)
features, label = next(iter(dataloaders["train"]))
print(features.keys(), label.shape)
\end{Verbatim}

\paragraph{Q: Full source code of the Python file}
\textbf{A:} The complete file is included verbatim for reference.

\VerbatimInput[fontsize=\scriptsize]{/workspace/dataloader/kuairec.py}

\end{CJK*}
\end{document}